% !PRINTED 537
\weberpage{537}
\begin{center}
\S.\,163.

\gesperrt{Reduction der Abel'schen Gleichungen auf cyklische.}
\end{center}

Wir haben schon vorhin gesehen, dass in einer transitiven \gesperrt{Abel'schen} Gruppe ausser der identischen keine Permutation vorkommt, die eine Ziffer ungeändert lässt. Daraus aber folgt noch ein anderer Satz. Nehmen wir an, eine Permutation $\pi$ einer solchen Gruppe sei in ihre Cyklen zerlegt und der Cyklus, der die wenigsten Glieder enthält, sei ein $r$-gliedriger. Dann wird $\pi^r$ die Glieder dieses Cyklus ungeändert lassen und muss also die identische Permutation sein. Daraus folgt aber, dass auch alle übrigen Cyklen von $\pi$, wenn noch andere vorhanden sind, aus $r$ Gliedern bestehen müssen, also:

\gesperrt{1. Eine Permutation einer transitiven Abel'schen Gruppe enthält nur Cyklen von gleicher Gliederzahl.}

Die Anzahl $r$ der Glieder eines Cyklus muss ein Theiler von $m$ sein, wenn $m$ der Grad der Gruppe ist, und wenn $m=rs$ ist, so ist $s$ die Anzahl der $r$-gliedrigen Cyklen, aus denen $\pi$ besteht.

Ist nun $P$ die Gruppe einer \gesperrt{Abel'schen} Gleichung $F(x)=0$, so nehmen wir irgend eine nicht identische Permutation $\pi$ aus $P$ heraus und zerlegen sie in ihre Cykeln
\weq{1}{\pi = \gamma\gamma_1\gamma_2\cdots\gamma_{s-1},}
worin jeder der Cyklen $\gamma$ sich auf $r$ Wurzeln der Gleichung $F(x)=0$ bezieht. Wir wollen diese Wurzeln so anordnen, dass
\weq{2}{\begin{aligned}
\gamma &= (\alpha,\alpha_1,\alpha_2\ldots\alpha_{r-1})\\
\gamma_1 &= (\beta,\beta_1,\beta_2\ldots\beta_{r-1})\\
&\ \,\vdots\\
\gamma_{s-1} &= (\sigma,\sigma_1,\sigma_2\ldots\sigma_{r-1})
\end{aligned}}
wird.

Ist $\pi_1$ irgend eine Permutation von $P$, so ist wegen der Vertauschbarkeit
\weq{3}{\pi_1^{-1}\pi\pi_1 = \pi.}

Nach dem Satze über die Bildung der Permutationen $\pi_1^{-1}\pi\pi_1$ (\S.\,154, 6.) darf also $\pi$ nicht geändert werden, wenn die Permutationen $\pi_1$ in den Cyklen von $\pi$ ausgeführt werden. Da aber die Cyklen vollständig bestimmt sind, abgesehen von ihrer An-
%CONT>>
